\section*{Discussion}

The core insight is that relevance is not applicability. Retrieval can return a
well-supported statement while a system still applies it to the wrong time,
scope, epistemic state or possible world. Situation engineering turns this
failure from an informal complaint about ``missing context'' into a set of
state variables, update rules and behavioural consequences that can be tested
separately.

\subsection*{A complementary layer, not a replacement vocabulary}

Prompt, context, retrieval, memory, tool, harness, evaluation, security and
serving engineering retain distinct responsibilities. Prompts specify the task;
context and retrieval supply evidence; memory preserves observations; tools
provide capabilities; harnesses govern workflow and permissions; evaluations
measure behaviour; serving systems manage resource constraints. Situation
engineering owns the query-relative semantics of the active state and the
propagation of unresolved applicability into policy. A useful contract is
\begin{equation}
(C_t,M_t,O_t,H_t)\rightarrow(\mathcal{S}_t,U_t,P_t),
\end{equation}
where context, memory, tool observations and harness traces produce candidate
situations $\mathcal{S}_t$, unresolved variables $U_t$ and the permitted policy
$P_t$.

This contract prevents conceptual inflation. Retrieval quality should not be
renamed situation engineering, and a JSON field named ``situation'' is not
evidence of situation awareness. A system must show that state representation
causally changes behaviour in predicted cases, that state uncertainty affects
answering or action, and that failures can be assigned to sensing, update,
validation, policy or generation.

\subsection*{Implications for hallucination and agent evaluation}

A generated claim can be textually entailed by a document yet wrong for the
query. Evaluation should therefore record both support and applicability. In
free-form generation, claims should be linked to the situation state and
provenance on which they depend. In agents, tools should expose preconditions,
state transitions and failure semantics rather than only free-text outputs;
harnesses should block consequential actions while answer-critical state is
unresolved. Evaluation engineering becomes the sensor layer that checks each
transition and not merely the final answer.

The response policy is also part of correctness. Clarification and abstention
are not failures when the situation is under-specified. Conversely, excessive
clarification is costly and can make systems unusable. Future evaluations should
therefore report answer accuracy, clarification utility, unnecessary-question
rate, abstention precision, selective risk, coverage, latency, token use and
monetary cost. These measures make explicit the trade-off between semantic
verification and serving efficiency.

\subsection*{What is established and what remains to be tested}

This study demonstrates a formal error class, an interoperable state schema and
category-specific causal effects of matched state ablations. It also shows that
performance degrades sharply when state must be recovered from partially
observable text. Together, these results identify situation sensing as a
separate engineering and evaluation target rather than treating every failure
as retrieval or generation error.

The controlled design cannot decide whether the full interface improves a
contemporary language model beyond a matched structured prompt or temporal RAG.
That comparison is the confirmatory submission gate for the general LLM claim:
same model, evidence and computation budget; independently authored natural
questions; direct, structured, temporal-retrieval, verification-agent and
situation conditions; and paired, scenario-clustered inference. The released
protocol makes this contrast executable, while the present results make no
claim from incomplete pilot calls. The SE0--SE4 maturity model and longer
research programme are provided in Supplementary Information.

\subsection*{Limitations and risks}

The diagnostic data are generated from a schema shared with the solver, so
benchmark-solver circularity is intentional and limits external validity. The
item count overstates statistical independence because examples share template
families. The baselines are software controls, not state-of-the-art LLM
systems. Confidence values are rule-derived and do not constitute neural
calibration. No human participants, multilingual evaluation, multimodal sensing or
deployed-agent study was completed. A small Gemini API pilot verified the
execution path but ended at unequal category coverage when free-tier quota was
exhausted; it is excluded from scientific results. These limitations
are not resolved by additional prose and require new experiments.

Explicit state can also create new attack surfaces. An adversary may manipulate
timestamps, source authority, identity, scope or world labels. A centralized
state representation may expose sensitive user information or create false
confidence in a flawed schema. Situation-engineered systems therefore require
provenance, access control, privacy-preserving memory, conflict visibility and
independent evaluation. ``Explicit'' must not be confused with ``correct''.

\subsection*{Conclusion}

Situation engineering provides a bounded and testable interface between
retrieval, state estimation and response policy. The controlled experiments
demonstrate that applicability variables have distinct behavioural consequences
and that imperfect sensing can erase the advantage of correct update rules.
Establishing broad utility now requires the pre-specified natural-text,
multi-model comparison against equally resourced structured prompting,
temporal retrieval and agent verification.
